\subsection{Post-saturation interoperability and provenance pressure}

The successor search further narrows the residual.  I-ADOPT already supplies a
cross-domain semantic decomposition for machine-interpretable scientific
variables~\cite{iadopt2021}; P3 therefore does not own structured variable
semantics.  FAIR~2.0 explicitly distinguishes terminological, ontological,
referential, schema and logical interoperability and argues that semantic
interoperability requires explicit mappings rather than a single privileged
ontology~\cite{fair2_2024}.  More recent provenance work goes further:
DEC interprets provenance-bearing statements as epistemic stance and preserves
source-relative disagreement without collapsing every attributed assertion into
a single factual core~\cite{dec2026}.  Consequently, neither provenance-aware
pluralism nor generic semantic interoperability is a P3 novelty claim.

The remaining P3 question is downstream and claim-relative: after a strong
semantic integration stack has produced structured candidate mappings, what
additional conditions license a \emph{scientific identity} assertion?  P3-X
therefore treats ontology/schema alignment, provenance, plural-view machinery
and consistency checks as donor inputs and tests only the joint decision
contract that returns \glue{}, \obstruction{}, \pluralview{} or \unresolved{}
from load-bearing construct, measurement, context, provenance-stance and
global-composition conditions.  An ideal product given exactly those same
coordinates and rules is expected to tie; the contribution is not an inherent
expressivity advantage.
